\section{Reading execution plans and SIMD dispatch}
\label{sec:appendix-explain}

\texttt{EXPLAIN} describes the selected plan. Adding \texttt{ANALYZE} executes
it and records work. Use text output to follow the tree and
\texttt{FORMAT JSON} to inspect named, nested properties. Some BioSeqDB nodes
print runtime property names even before execution; a zero or
\texttt{not executed} value in plain \texttt{EXPLAIN} is not a measurement.
The examples below illustrate plan inspection. They are not benchmark results.

\subsection{Start with the query and its PostgreSQL parent nodes}

The following example keeps its table and settings inside a transaction.
It projects match coordinates and a CIGAR string at the search level, where
those values belong to the current match. The final rollback removes the
example table and restores the local settings.

\par\noindent\begin{minipage}{\linewidth}
\begin{lstlisting}[style=bioseq-appendix-sql]
BEGIN;
SET LOCAL max_parallel_workers_per_gather = 0;
SET LOCAL bioseq.execution.parallel.enabled = off;
SET LOCAL bioseq.index.insert_maintenance.enabled = on;
SET LOCAL bioseq.index.sequence_read.enabled = off;
SET LOCAL bioseq.vectorization.enabled = on;
SET LOCAL bioseq.aligner.vulkan.enabled = off;

CREATE TABLE explain_reads (
  pk integer PRIMARY KEY,
  seq bioseq.seq('dna4', 32) NOT NULL
) USING bioseq;
INSERT INTO explain_reads VALUES
  (1, 'TCACTGTAGGATCGTAGCTAGTCCGTACTGTAGCTGATCGTAGCCCTAG'
      'TCACTGTAGCCTGATCGTAGCTAGTCACTGTAGCCATACCTGAA'),
  (2, 'CCCCAACTTAGTCACTTCTCAAAAATTACTCTCTCATCATCATCTACTC'
      'CTAGTCACTGTAGCCTGATCGTAGCTAGTCACACACACCCCAAA');

EXPLAIN (ANALYZE, VERBOSE, BUFFERS, COSTS OFF)
SELECT pk,
       bioseq.target_start(pk) AS target_start,
       bioseq.target_end(pk) AS target_end,
       bioseq.cigar(pk) AS cigar
FROM explain_reads
WHERE public.mmsearch(
  seq, 'CTAGTCACTGTAGCCTGATCGTAGCTAGT'::bioseq.seq('dna4', 32))
ORDER BY pk, target_start;
ROLLBACK;
\end{lstlisting}
\end{minipage}

\begin{itemize}
  \item \textbf{Before running.} Use \texttt{EXPLAIN (VERBOSE, COSTS OFF)}
    with the same query to inspect its shape without running the search.
  \item \textbf{For structured output.} Add \texttt{FORMAT JSON} to either
    form. In JSON, identify a custom node using both \texttt{Node Type} and
    \texttt{Custom Plan Provider}.
  \item \textbf{For application parameters.} A value such as
    \texttt{\$1::bioseq.seq('dna4', 32)} requires a prepared statement or a
    parameter supplied by the client. Use a typed literal in a standalone
    statement.
  \item \textbf{For expensive work.} \texttt{EXPLAIN ANALYZE} performs the
    alignments, remote requests and cache waits required by the statement.
\end{itemize}

\begingroup\small
\begin{longtable}{@{}p{.35\linewidth}p{.60\linewidth}@{}}
\toprule
\textbf{PostgreSQL property} & \textbf{How to read it}\\
\midrule\endhead
\texttt{actual time}, \texttt{rows}, \texttt{loops} &
Times and rows are averages per loop. Multiply by the loop count when
estimating total work at that node. A parent includes work in its children;
do not sum all node times.\\
\texttt{Buffers} & PostgreSQL shared, local or temporary blocks touched.
These are not sequence symbols, packed bytes or GPU upload bytes.\\
\texttt{Rows Removed by Filter} & Rows rejected by the ordinary filter
printed for this node.\\
\texttt{Sort Method}, \texttt{Memory} & How a PostgreSQL parent sort ran,
including whether it required disk work.\\
\texttt{Planning Time}, \texttt{Execution Time} & Whole-statement planning
and execution times. Compare these with the narrower BioSeqDB counters.\\
\bottomrule
\end{longtable}
\endgroup

A \texttt{Sort}, \texttt{Aggregate}, \texttt{Materialize}, \texttt{Limit}, or
join above the search can dominate the statement. For example,
\texttt{ORDER BY bioseq.score(pk) DESC, pk} sorts match scores using a
PostgreSQL sort; it does not request the sequence batch sorter.

\subsection{Identify the custom provider}

\begingroup\small
\begin{longtable}{@{}p{.35\linewidth}p{.60\linewidth}@{}}
\toprule
\textbf{Custom Plan Provider} & \textbf{Responsibility}\\
\midrule\endhead
\texttt{BioseqScan} & Read a local BioSeqDB table and, when owned by this
node, evaluate a sequence search.\\
\texttt{VectorAggNode} & Apply an eligible aggregate directly to column
batches from its \texttt{BioseqScan} child.\\
\texttt{BioseqIndexScan} & Use a PostgreSQL index with BioSeqDB tuple fetch.\\
\texttt{BioseqSliceProject} & Evaluate a lazy sequence expression and
materialize values at required boundaries.\\
\texttt{BioseqBatchJoin} & Own a sequence join, or a search over the
per-row expressions produced by a child join.\\
\texttt{BioseqBatchSortScan} & Buffer and order one eligible sequence key.\\
\texttt{BioseqCacheScan} & Produce, replay, reuse or bypass query results.\\
\texttt{BioseqExternalScan} & Read a remote provider-backed relation.\\
\bottomrule
\end{longtable}
\endgroup

For example, the following is a \emph{schematic} tree. It shows responsibility
and contains no measured costs or times.

\par\noindent\begin{minipage}{\linewidth}
\begin{lstlisting}[style=bioseq-appendix-text]
Sort                         -- order the final match rows
  Custom Scan (BioseqScan)    -- read targets and produce matches
\end{lstlisting}
\end{minipage}

\subsection{Local scan: columns, filters and index sources}

The \texttt{Bioseq Projected Columns} property lists what the parent needs.
Check it first: a query that needs only metadata should not return a whole
sequence merely because an application normally includes that column.

\begingroup\small
\begin{longtable}{@{}p{.46\linewidth}p{.49\linewidth}@{}}
\toprule
\textbf{BioseqScan property} & \textbf{Meaning}\\
\midrule\endhead
\texttt{Bioseq Chunk Group Filters} & Qualifiers considered for pruning
whole chunk groups.\\
\texttt{Bioseq Chunk Groups Removed by Filter} & Groups removed by the
pruning path, when scan statistics are available.\\
\texttt{Bioseq Vectorized Filter} & Qualifiers assigned to the vectorized
filter evaluator.\\
\texttt{Bioseq Fetched From Meta} & Projections answered from metadata,
including a requested symbol and its alphabet rank for symbol counts.\\
\texttt{Bioseq Index Read Requirement} & The target-index source rule:
a route may permit, require or exclude stored target-index data.\\
\texttt{Bioseq Match Operand Evaluation} & A target or query is evaluated
as an expression. An arbitrary transformed target cannot simply use the
stored index of the original column.\\
\texttt{Bioseq Typmod Coalesce} & A compatible query representation is
re-encoded to the target alphabet and active storage width. It does not
translate between different biological alphabets.\\
\texttt{Bioseq Fetched From Sequence Index} & Plan-time estimate of rows
with suitable stored target-index data.\\
\texttt{Bioseq Recomputed} & Plan-time estimate of rows whose minimizers
must be rebuilt, for example when index reading is off or data is absent.\\
\texttt{Bioseq Scheduler} & A point-in-time scheduler snapshot: supervisor
state, concurrency, workers, batches, waiting work, query rates, asynchronous
persistence and reserved index memory. It is not query-stage timing.\\
\bottomrule
\end{longtable}
\endgroup

For a metadata query, compare these two projections:

\par\noindent\begin{minipage}{\linewidth}
\begin{lstlisting}[style=bioseq-appendix-sql]
EXPLAIN (ANALYZE, VERBOSE, COSTS OFF)
SELECT pk, bioseq.symbol_count(seq, 'A') FROM reads;

EXPLAIN (ANALYZE, VERBOSE, COSTS OFF)
SELECT pk, bioseq.symbol_count(seq, 'A'), seq::text FROM reads;
\end{lstlisting}
\end{minipage}

Here \texttt{reads} is an existing table with \texttt{pk} and \texttt{seq}
columns. Metadata can answer the count in both statements; the second also
requires sequence payloads for text output. Compare estimated index shares
with the actual row counts below, rather than treating a percentage in the
plan as proof that every row used the stored index.

\subsection{Pipeline and Highway fields}
\label{sec:appendix-highway}

In JSON, \texttt{Bioseq Pipeline} contains the algorithm description and its
stages. Text output presents a tree ending in \texttt{Submit results}.
The top-level properties are \texttt{Chaining Algorithm},
\texttt{Base Alignment Algorithm}, \texttt{Pipeline Concurrency Mode},
\texttt{Max Concurrency}, and \texttt{Index Memory Reserved Bytes}.
The algorithm names describe the parameters captured for this execution.
Concurrency and reserved-memory values describe scheduler state; they do not
mean every available thread worked on every stage.

Highway chooses one CPU instruction-set target from the compiled targets
that the host can run. The processing entry binds and checks its kernel table
once per batch. It does not choose the fastest biological algorithm by
benchmarking alternatives. A target name can change with the extension build,
CPU or operating-system support.

\begingroup\small
\begin{longtable}{@{}p{.46\linewidth}p{.49\linewidth}@{}}
\toprule
\textbf{Highway SIMD Dispatch property} & \textbf{Meaning}\\
\midrule\endhead
\texttt{Highway Compiled Targets} & Kernels included in this binary.\\
\texttt{CPU-Supported Targets} & Targets allowed by the host CPU and
operating system.\\
\texttt{Runnable Targets} & The intersection of the preceding sets.\\
\texttt{Planned Architecture} & One selected runnable target.\\
\texttt{Executed Architecture} & Target recorded by execution. Analyzed output can omit it
when no target ran. \texttt{mixed} means the combined records contain
more than one executed target.\\
\bottomrule
\end{longtable}
\endgroup

\begin{figure}[htbp]
\centering
\begin{tikzpicture}[font=\small,x=1cm,y=1cm]
  \draw[fill=appendixBlue!12,draw=appendixBlue!85!black]
    (3.20,1.60) ellipse (2.50 and 1.28);
  \draw[fill=appendixOrange!12,fill opacity=.65,
        draw=appendixOrange!90!black]
    (6.20,1.60) ellipse (2.50 and 1.28);
  \node[align=center] at (1.92,1.75) {compiled\\targets};
  \node[align=center] at (7.45,1.75) {CPU/OS\\supported};
  \node[align=center,font=\footnotesize] at (4.70,1.86)
    {runnable\\targets};
  \fill[appendixGreen!80!black] (4.70,1.06) circle (2.8pt);
  \draw[-{Latex[length=1.7mm]},thick]
    (4.70,.97) -- (4.70,-.08);
  \node[anchor=north,align=center] at (4.70,-.18)
    {one selected target\\checked when the batch binds its kernels};
\end{tikzpicture}
\caption{A compiled SIMD kernel is usable only if the host supports it.
The intersection contains the runnable targets. The dot denotes the
selected target.}
\label{fig:appendix-highway-sets}
\end{figure}

\texttt{Planned Highway Functions} lists top-level planned functions.
Each \texttt{Highway Function} object has \texttt{Function},
\texttt{Planned Architecture} and, at the top level, \texttt{Stage}.
Stage-owned entries instead appear under that stage's
\texttt{Highway Functions}; their surrounding stage already identifies the
owner. A planned function entry is not proof of execution.

Each object in \texttt{Pipeline Stages} can contain
\texttt{Stage Name}, \texttt{Algorithm}, algorithm-specific parameters,
computed properties, \texttt{Highway Functions}, \texttt{Description}, and
\texttt{Estimated Batch Size}. The description substitutes for an algorithm
object when appropriate. Disabled stages are absent. CIGAR, divergence,
mapping-quality and backend choices can therefore change the list.

For a CPU dispatch comparison, execute the same search with these settings:

\par\noindent\begin{minipage}{\linewidth}
\begin{lstlisting}[style=bioseq-appendix-sql]
BEGIN;
SET LOCAL bioseq.aligner.vulkan.enabled = off;
SET LOCAL bioseq.aligner.minimizer.algorithm = 'two_stack_nthash';
SET LOCAL bioseq.aligner.chaining.algorithm = 'simd_minimap2';
SET LOCAL bioseq.aligner.base_alignment.algorithm = 'block_aligner';
SET LOCAL bioseq.vectorization.enabled = on;
-- Run the search with EXPLAIN (ANALYZE, VERBOSE, FORMAT JSON).
SET LOCAL bioseq.vectorization.enabled = off;
-- Repeat the same query, then compare its result rows and both plans.
ROLLBACK;
\end{lstlisting}
\end{minipage}

The three algorithm settings select the minimizer, chainer and base aligner.
\nolinkurl{bioseq.vectorization.enabled} permits vectorized routes; it does not
name a CPU instruction set or guarantee that every Highway-backed helper disappears when
off. Read the per-stage \texttt{vectorised} labels and the whole dispatch
object. Check that the planned target belongs to the runnable set, that
planned functions agree with it, and that analyzed execution reports the
expected target. An incomplete table, no runnable target or a binding mismatch
raises an error for the selected SIMD work. Keep the full target inventory
with performance comparisons.

\subsection{Observed search work}

\texttt{Bioseq Execution Statistics} groups the search counters in JSON.
Text output uses compact backend, index, GPU, host and prefilter lines.
Read these counters from \texttt{EXPLAIN ANALYZE}; their names can otherwise
appear with unexecuted values on some nodes.

\begingroup\small
\begin{longtable}{@{}p{.49\linewidth}p{.46\linewidth}@{}}
\toprule
\textbf{Property} & \textbf{Meaning}\\
\midrule\endhead
\texttt{Requested Backend}, \texttt{Actual Backend} & Requested and
completed processing routes.\\
\texttt{Requested Target Index Source},
\texttt{Actual Target Index Source} & Requested and used target-index
sources. Compare these before interpreting index settings.\\
\texttt{Target Rows} & Target rows processed.\\
\texttt{Persisted Index Rows} & Rows served by stored target-index data.\\
\texttt{Missing Persisted Index Rows} & Rows for which expected stored
data was missing.\\
\texttt{Recomputed Target Index Rows} & Rows whose target index was rebuilt.\\
\texttt{GPU Seed Dispatches}, \texttt{GPU Chain Dispatches},
\texttt{GPU Chain Backtrack Dispatches} & Device dispatch counts for
seeding, forward chaining and backtracking.\\
\texttt{GPU Coordinate Vote Dispatches}, \texttt{GPU Extend Dispatches}
& Dispatch counts for coordinate voting and extension.\\
\texttt{Backend Declines}, \texttt{Backend Fallbacks} & Requested routes
declined, and work that returned to another backend. These counters do not
promise that every unsupported GPU request has a CPU fallback.\\
\texttt{Query Sketch Host Wall Ns},
\texttt{Target Index Build Host Wall Ns} & Host wall time in nanoseconds
for the named work; not PostgreSQL whole-node time.\\
\texttt{Host Query Hash Readbacks} & Query-hash readbacks to the host.\\
\texttt{Host Query Symbol Packing Calls},
\texttt{Host Target Symbol Packing Calls} & Calls packing query or target
symbols on the host.\\
\texttt{Host Alignment Job Classification Calls},
\texttt{Host CIGAR Construction Calls} & Host job-classification and CIGAR
construction calls.\\
\bottomrule
\end{longtable}
\endgroup

A nonzero stage dispatch counter shows that work reached that device stage.
It does not mean the complete statement ran on the GPU, or that host work was
zero. The text \texttt{Submit results} line can also include target and
index-source row counts.

\subsubsection{Prefilter counters}

The prefilter rejects candidate work before full alignment. When statistics
were collected, \texttt{Prefilter Stats Available} is true and the remaining
fields identify the work. The implementation omits this group when statistics
are unavailable; it need not print a false-valued availability property.
All names in the following table have the prefix \texttt{Prefilter }.

\begingroup\small
\begin{longtable}{@{}p{.49\linewidth}p{.46\linewidth}@{}}
\toprule
\textbf{Property suffix} & \textbf{Meaning}\\
\midrule\endhead
\texttt{Resolved Algorithm}, \texttt{Resolved Backend} & Algorithm and
backend selected after route checks.\\
\texttt{Stats Saturated}, \texttt{Stats Consistent} & Whether a counter
reached its numeric limit, and whether the counters satisfy their checked
relationships. Neither is a claim about biological sensitivity.\\
\texttt{Dispatches}, \texttt{SneakySnake Dispatches},
\texttt{GateKeeper Dispatches} & Total dispatches and counts by algorithm.\\
\texttt{Scalar Dispatches}, \texttt{SIMD Dispatches},
\texttt{Vulkan Dispatches} & Counts by backend. These classify the same
work along another axis; do not add them to algorithm counts.\\
\texttt{Jobs}, \texttt{Evaluated Jobs}, \texttt{Length 64 Jobs},
\texttt{Unequal-Length Jobs}, \texttt{Unsupported Jobs} & Workload counts:
submitted, evaluated, length-64, unequal-length and unsupported jobs.\\
\texttt{Rejected Jobs}, \texttt{Rejected Regions},
\texttt{Surviving Regions} & Outcomes at job and candidate-region level.\\
\texttt{Mask Rejection Candidates}, \texttt{Exact Guard Activations},
\texttt{Guard Confirmed Rejections}, \texttt{Hamming Exact Rejections}
& Candidates flagged by a rejection mask, exact guard checks, rejections
confirmed by that guard, and exact Hamming-based rejections.\\
\bottomrule
\end{longtable}
\endgroup

A high rejection rate is useful only when it saves more expensive work while
meeting the required sensitivity. Compare surviving work and total query time,
not just the number rejected.

\subsubsection{Sequence column cache}

With \nolinkurl{bioseq.column_cache.enabled}, a scan can print a \texttt{Cache}
section with \texttt{Cache Hits}, \texttt{Cache Misses},
\texttt{Cache Evictions}, \texttt{Cache Writes},
\texttt{Cache Maximum Size}, \texttt{Cache Ending Size}, and
\texttt{Total Cache Entries}. These are column-cache snapshots from the
backend's cache statistics, not isolated deltas for this query.
\texttt{Cache Maximum Size} is the observed peak payload allocation, not a
configured limit; \texttt{Cache Ending Size} is the current payload allocation.
This cache
holds sequence-column data; the query-result cache has its own
\texttt{BioseqCacheScan} provider and action fields.

\subsection{Lazy slices and materialization}

\texttt{BioseqSliceProject} reports \texttt{Projection Mode}
(\texttt{source-backed lazy}), \texttt{Source Pieces},
\texttt{Estimated Piece Boundaries}, and an optional \texttt{Fallback Reason}.
A piece is a source range in the lazy expression. The following pair separates
length calculation from returning concrete sequence text:

\par\noindent\begin{minipage}{\linewidth}
\begin{lstlisting}[style=bioseq-appendix-sql]
EXPLAIN (ANALYZE, VERBOSE, COSTS OFF, FORMAT JSON)
SELECT bioseq.len(seq[10:-10]) FROM reads;

EXPLAIN (ANALYZE, VERBOSE, COSTS OFF, FORMAT JSON)
SELECT seq[10:-10]::text FROM reads;
\end{lstlisting}
\end{minipage}

The \texttt{Slice Recipe} array contains \texttt{Recipe Node} objects with
\texttt{Stable ID}, \texttt{Opcode}, \texttt{Expression},
\texttt{Depends On}, and applicable \texttt{Lower Bound} or
\texttt{Upper Bound}. Dependencies refer to earlier recipe values.
The operations cover source, literal, element, slice, concatenation, reverse,
reverse complement, length, cardinality, literal bridge and materialization.
An omitted bound remains omitted; it is not a stored numeric bound.
A recipe can transform source ranges without copying a complete sequence.

\begingroup\small
\begin{longtable}{@{}p{.46\linewidth}p{.49\linewidth}@{}}
\toprule
\textbf{Slice Materialization property} & \textbf{Meaning}\\
\midrule\endhead
\texttt{Variable}, \texttt{Expression} & Boundary value and SQL expression
that needs a concrete representation.\\
\texttt{Mechanism}, \texttt{Reason} & How and why it is materialized.\\
\texttt{Depends On}, \texttt{Encoding}, \texttt{Cache Scope} & Earlier
boundaries, result encoding and reuse scope.\\
\texttt{Producer Value}, \texttt{Source Provenance} & Recipe producer and
source expression, included with \texttt{VERBOSE}.\\
\texttt{Estimated Symbols}, \texttt{Estimated Bytes} & Estimates when
costs are shown. The current byte estimate copies the estimated symbol count;
it is a rough one-byte-per-symbol estimate, not predicted packed storage.\\
\texttt{Actual Calls}, \texttt{Actual Input Symbols},
\texttt{Actual Output Symbols}, \texttt{Actual Bytes} & Recorded calls,
input/output symbol counts and concrete bytes from analyzed execution.\\
\texttt{Reused Calls}, \texttt{Peak Retained Bytes} & Calls served by reuse
and peak retained boundary bytes.\\
\texttt{Counter Overflow}, \texttt{Incomplete Workers} & Counter overflow
or incomplete collection of parallel-worker records.\\
\texttt{Actual Time} & Recorded boundary time when timing is enabled.\\
\bottomrule
\end{longtable}
\endgroup

Text output with costs also prints \texttt{Planned Materialized Symbols} and
\texttt{Planned Materialized Bytes}. An empty materialization list means no
concrete payload boundary was recorded for that recipe. It does not describe
unrelated allocations elsewhere in the statement.

\subsection{Sequence joins}

For each \texttt{BioseqBatchJoin}, first identify its target and query aliases.
\texttt{Bioseq Residual Mmsearch Qual: absent} means the custom node owns
the search predicate. An ordinary residual search predicate is a reason to
check whether the intended rewrite occurred.

\begingroup\small
\begin{longtable}{@{}p{.48\linewidth}p{.47\linewidth}@{}}
\toprule
\textbf{BioseqBatchJoin property} & \textbf{Meaning}\\
\midrule\endhead
\texttt{Bioseq Query Batch Capacity} & Capacity of the runtime query buffer.\\
\texttt{Bioseq Index Read Requirement} & Target-index source rule.\\
\texttt{Bioseq Query Planner Rows}, \texttt{Bioseq Query Planner Width}
& Estimated query rows and bytes per row.\\
\texttt{Bioseq Query Planner Admission} & Planned query materialization
bytes.\\
\texttt{Bioseq Query Actual Materialization} & Query materialization bytes
observed during execution.\\
\texttt{Bioseq Distinct Targets}, \texttt{Bioseq Query Batches} & Distinct
non-NULL target child tuples processed and executed query batches. Despite
the name, the target counter does not independently apply SQL DISTINCT.\\
\bottomrule
\end{longtable}
\endgroup

Additional join diagnostics are useful when a requested shape is rejected:
\begin{itemize}
  \item \texttt{Bioseq CPU Prepared Target} describes target preparation:
    a disabled route, a grouped-SoA reader retained across query waves, or a
    query-hash union with one stored target index per parameter group.
    The grouped-SoA label covers a wider stored-index route: a compatible
    hash-major reader can be selected, and the whole-sequence reader loads
    relation payload. The label alone therefore does not prove that a mapped
    file was read or that CPU or GPU work ran.
  \item \texttt{Bioseq Honest Startup Cost},
    \texttt{Bioseq Honest Total Cost}, and
    \texttt{Bioseq Incumbent Cost Floor} expose the custom route's estimated
    costs and comparison floor. Their units are planner cost units, not time.
  \item \texttt{Bioseq Batch ABI Version} identifies the serialized
    plan format.
  \item \texttt{Bioseq Match Placeholder Sites} counts places filled
    with current-match values.
  \item \texttt{Bioseq Requirement Executable} reports whether the decoded
    index-source requirement was accepted.
\end{itemize}

These properties can appear before work runs. Use analyzed values when
comparing planned and actual materialization. The node also prints the shared
Highway and execution-statistics sections. Read each join separately when the
statement has more than one sequence search.

A transformed target can instead produce this \emph{schematic} ownership
header under the same provider:

\par\noindent\begin{minipage}{\linewidth}
\begin{lstlisting}[style=bioseq-appendix-text]
Custom Scan (BioseqBatchJoin)
  Bioseq Mmsearch Ownership: per-row recipe
  Bioseq Reconstructed Join: child plan
  Bioseq Residual Mmsearch Qual: absent
\end{lstlisting}
\end{minipage}

This form evaluates search operands for each row supplied by its child plan.
Its \texttt{Bioseq Batch Capacity} describes its buffer;
\texttt{Bioseq Index Read Requirement} is absent for an arbitrary transformed
target because the original column's stored index does not describe that
logical value. Runtime initialization adds the shared dispatch and work
statistics. The additional properties
\texttt{Bioseq Owned RestrictInfo Path Copies},
\texttt{Bioseq Executable Mmsearch Qual Copies},
\texttt{Bioseq Vector Mmsearch Qual Copies},
\texttt{Bioseq Storage Mmsearch Qual Copies}, and
\texttt{Bioseq Descendant Mmsearch Owner Copies} count copies of the owned
predicate found in different parts of the retained plan. They help detect a
predicate that would be evaluated again outside its intended owner. They are
plan diagnostics, not counts of sequence matches.

An ordinary predicate written next to a join does not prove that it reduced
query or target admission. Read both children, actual materialization,
distinct targets and the final result boundary. When the boundary must be
explicit, stage filtered input in BioSeqDB tables or filter already
materialized match rows. Coordinate helpers should remain plain projections
at the search level; add slice or provider offsets in an outer query over the
materialized match columns.

\subsection{Vulkan availability and sequence sorting}

\texttt{bioseq.aligner.vulkan.enabled} permits eligible Vulkan work. Inspect
the runtime separately from that preference:

\par\noindent\begin{minipage}{\linewidth}
\begin{lstlisting}[style=bioseq-appendix-sql]
SELECT is_available, device_name, max_workgroup_count_x,
       has_scalar_block_layout, device_local_bytes,
       is_software_device, initialization_error
FROM bioseq.vulkan_runtime_status();
\end{lstlisting}
\end{minipage}

\texttt{is\_available} is the first gate. If false, inspect
\texttt{initialization\_error}. Otherwise retain the device name, memory and
capabilities with any measured plan. A software Vulkan device is distinct
from a physical GPU. The remaining fields have specific units:
\begin{itemize}
  \item \texttt{max\_workgroup\_count\_x}: maximum workgroup count in
    the first dispatch dimension;
  \item \texttt{has\_scalar\_block\_layout}: support for the scalar
    buffer-layout feature used by eligible kernels;
  \item \texttt{device\_local\_bytes}: reported device-local memory in
    bytes, not memory already reserved for this query.
\end{itemize}

The supported Vulkan alignment route uses five more settings under the
\nolinkurl{bioseq.aligner.vulkan.} prefix:
\begin{itemize}
  \item \nolinkurl{seed_lookup.enabled};
  \item \nolinkurl{index_cache.enabled};
  \item \nolinkurl{fused_frontend.enabled};
  \item \nolinkurl{index_build.enabled};
  \item \nolinkurl{alignment_extension.enabled}.
\end{itemize}
Its minimizer and chainer values are \nolinkurl{vulkan_minimap_integer} and
\nolinkurl{vulkan_tiled_dynamic_programming}. These settings form a coordinated route
with representation, device and query-shape requirements, not independent
promises that any arbitrary search runs on the GPU. The stored grouped
structure-of-arrays route currently returns coordinates; it does not provide
an arbitrary CIGAR route. Highway fields retain their CPU dispatch meaning.

The separate planner switch
\nolinkurl{bioseq.planner.sequence_sort_scan.enabled}, off by default,
permits a sequence batch-sort custom node. Eligibility requires one direct
sequence ordering key and a supported child path.

\begingroup\small
\begin{longtable}{@{}p{.41\linewidth}p{.54\linewidth}@{}}
\toprule
\textbf{Query shape} & \textbf{Planning consequence}\\
\midrule\endhead
\texttt{ORDER BY seq} & Can use \texttt{BioseqBatchSortScan}.\\
\texttt{ORDER BY seq DESC} & Can use the custom sorter with descending
sequence order.\\
\texttt{ORDER BY seq, pk} & More than one key: retain an ordinary sort.\\
\texttt{ORDER BY seq::text} & This is text ordering, not sequence batch
sorting.\\
Planner switch off & Keep the ordinary PostgreSQL sort.\\
\bottomrule
\end{longtable}
\endgroup

The result need not project the sequence key. For example, this comparison
returns only primary keys:

\par\noindent\begin{minipage}{\linewidth}
\begin{lstlisting}[style=bioseq-appendix-sql]
BEGIN;
CREATE TABLE explain_sort (
  pk integer PRIMARY KEY, seq bioseq.seq('dna4', 32)
) USING bioseq;
INSERT INTO explain_sort VALUES
  (1, 'TTTT'), (2, 'AAAA'), (3, 'CCCC'), (4, 'GGGG');
SET LOCAL bioseq.planner.sequence_sort_scan.enabled = on;
SET LOCAL bioseq.aligner.vulkan.enabled = off;
EXPLAIN (ANALYZE, VERBOSE)
SELECT pk FROM explain_sort ORDER BY seq;

SET LOCAL bioseq.aligner.vulkan.enabled = on;
EXPLAIN (ANALYZE, VERBOSE)
SELECT pk FROM explain_sort ORDER BY seq;
ROLLBACK;
\end{lstlisting}
\end{minipage}

A four-row fixture demonstrates the controls; it is too small to establish a
performance advantage. For a result comparison, run the same ordered
\texttt{SELECT} with the planner switch off and on before the rollback.
The four distinct keys must appear in the same order. Then inspect whether
Vulkan actually ran.

\begingroup\small
\begin{longtable}{@{}p{.49\linewidth}p{.46\linewidth}@{}}
\toprule
\textbf{BioseqBatchSortScan property} & \textbf{Meaning}\\
\midrule\endhead
\texttt{Bioseq Sort Accelerator} & Recorded route: \texttt{vulkan},
\texttt{stdc\_compare}, or
\texttt{stdc\_compare backend source}. Before execution it may be
uninitialized or \texttt{not executed}.\\
\texttt{Bioseq Sort Rows}, \texttt{Bioseq Sort Runtime Ns} & Batch rows
and sort wall time in nanoseconds. This timer excludes preceding source
admission and materialization.\\
\texttt{Bioseq GPU Sort Dispatches} & Sequence-sort device dispatches.\\
\texttt{Bioseq GPU Global Compare Dispatches} & Full-payload device
comparison dispatches.\\
\texttt{Bioseq GPU Dense Scatter Dispatches} & Prefix dense-scatter
iteration dispatches.\\
\texttt{Bioseq GPU Sort Uploaded Bytes} & Bytes uploaded by the recorded
GPU sort attempt.\\
\texttt{Bioseq Source Rows Admitted}, \texttt{Bioseq Source Bytes Admitted}
& Source rows and bytes admitted for the sort representation.\\
\texttt{Bioseq Source Admission Fallbacks} & Source-admission failures
that used the original source views.\\
\texttt{Bioseq Source Routes Acquired} & Whether the sorter acquired the
required source-access routes.\\
\bottomrule
\end{longtable}
\endgroup

\begin{itemize}
  \item \textbf{CPU route.} With Vulkan disabled, the sorter uses CPU
    comparison and reports zero GPU dispatches.
  \item \textbf{GPU route.} It needs a usable device, compatible packed
    input, bounded sizes, admitted memory and a successful device operation.
    DNA4 and DNA5 have dedicated kernel kinds; supported other alphabets use
    the generic kind.
  \item \textbf{Fallback.} Device, layout, memory, overflow or device-sort
    failure returns to CPU comparison. A source-admission failure increments
    its own counter and uses the \texttt{backend source} label.
  \item \textbf{Ordering.} The sorter may try a short prefix when sampled
    leading units vary. It accepts that route only with the required
    dense-scatter result; otherwise it retries with the complete payload for
    exact comparison. Descending order reverses the permutation. Equal
    sequences use a deterministic source-index tie break within this sorter;
    SQL does not promise a common tie order across unrelated plans.
\end{itemize}

To confirm device sorting, check two properties in analyzed output:
\begin{itemize}
  \item \texttt{Bioseq Sort Accelerator} must be \texttt{vulkan};
  \item \texttt{Bioseq GPU Sort Dispatches} must be greater than zero.
\end{itemize}
The node name or enabled GUC alone does not establish device sorting.
Sort counters reset on rescan;
do not automatically multiply them by PostgreSQL's loop count. Unlike this optional
sort accelerator, a requested strict Vulkan alignment route can reject an
incomplete route with an error.

\subsection{Query-cache nodes}

The query cache is experimental and insufficiently tested; a hit label does
not establish correct reuse under every supported plan and transaction.
Its operational examples appear in Section~\ref{sec:appendix-query-cache}.
Read the action, branches and residual work before interpreting timing.

\subsubsection{Exact final results}

For \texttt{Cache Mode: exact\_query},
\texttt{Cache Entry Kind} is \texttt{ExactFinalResult}. The plan reports
\texttt{Planned Cache Action}, \texttt{Cache Force Bypass},
\texttt{Cache Bypass Reason} when forced,
\texttt{Provider Runtime Recipes}, \texttt{Cache Parameter Count},
\texttt{Cache Dependency Count}, \texttt{Cache Dependencies}, and
\texttt{SQL}. With analyzed output it adds
\texttt{Provider Runtime Recipe State} and \texttt{Actual Cache Action}.
The actual action distinguishes:
\begin{itemize}
  \item \texttt{[Hit] exact final result}: replay an exact entry;
  \item \texttt{[Produce] exact final result}: execute and produce an entry;
  \item \texttt{[Bypass] retained child}: run the retained child plan.
\end{itemize}

\texttt{Cache Readiness}, \texttt{Cache Storage},
\texttt{Cache Storage Ownership}, and \texttt{Cache Janitor} are readable
status strings. They report readiness; storage attachment, reservations,
ceilings, manifests and waiters; readers and retained ownership references;
and the janitor lifecycle. They are not stable numeric tables. Use
\texttt{bioseq\_cache.stats()} and \texttt{bioseq\_cache.entries()} for tabular
inspection, with the monitoring limitations described in the cache tutorial.

BioSeqDB's provider detail can also show \texttt{Bioseq Cache Provider},
\texttt{Bioseq Provider Recipe Revision}, \texttt{Bioseq Provider Relation}
for a scan, or \texttt{Bioseq Provider Plan Node} for a batch join.
\texttt{Bioseq Exact Recipe Valid} reports recipe validation;
\texttt{Bioseq Reusable Coverage} is false for the current BioSeqDB adapters.

\subsubsection{Reusable coverage: generic provider interface}

The cache core implements the following coverage-plan fields, but the current
BioSeqDB adapters supply exact-result callbacks only. Selecting
\texttt{reusable\_fragments} does not add the missing coverage callbacks.
This reference describes what a provider implementing that interface can
report; it is not a claim that current BioSeqDB scans reuse partial ranges.

For \texttt{Cache Mode: reusable\_fragments},
\texttt{Cache Entry Kind} is \texttt{ReusableCoverage}. The top level reports
\texttt{Cache Force Bypass}, \texttt{Planned Cache Action},
\texttt{Actual Cache Action} after analysis, \texttt{Cache Entry},
\texttt{SQL}, \texttt{Reusable Sources}, \texttt{Selected Cache Branches},
and \texttt{Provider Runtime Recipe Valid}.

A current generation snapshot adds \texttt{State}, \texttt{Producer PID}
when permitted, \texttt{Age}, \texttt{Effective TTL}, \texttt{Rows},
\texttt{Logical Bytes}, \texttt{Stored Bytes}, and \texttt{Dependencies}.
Without one, state is \texttt{planned/provisional}; execution can still supply
its own rows and byte counts. The displayed \texttt{Age} is milliseconds
since the producer started. Despite its label, \texttt{Effective TTL}
currently shows the producer's TTL in milliseconds. Reuse checks elapsed time
since completion against the smaller producer and consumer TTL. Therefore the
two displayed numbers alone do not decide whether reuse is still allowed.

Each \texttt{Cache Branch} reports \texttt{Action} and, when needed,
\texttt{Recheck}. The \texttt{Residual Child} contains a \texttt{Predicate}
for uncached work. An optional \texttt{Provider} section supplies
provider-specific reusable-domain details. A query can reuse branches and
compute the remaining rows in the same execution.

\subsection{External scans, ordinary indexes and aggregates}

\texttt{BioseqExternalScan} identifies its logical relation using
\texttt{Bioseq External Entity}, and its adapter using
\texttt{Bioseq External Provider}. Read the ordinary PostgreSQL
\texttt{Filter} for provider and record restrictions. Plain
\texttt{EXPLAIN} does not perform the remote request;
\texttt{EXPLAIN ANALYZE} does. Actual rows and node timing are the main
execution measurements. A physical sequential scan is not a usable fallback
for these provider relations: their table access method rejects it.

\texttt{BioseqIndexScan} starts with \texttt{BioseqIndexScan using } and an
index name, followed by ordinary \texttt{Index Cond}, \texttt{Order By} and
\texttt{Filter} properties when applicable. Analyzed output reports
\texttt{Rows Removed by Index Recheck} and \texttt{Rows Removed by Filter};
text suppresses zero counts while structured output can retain them.
This node changes tuple fetching, not the meaning of a metadata index into a
sequence-alignment index.

\texttt{VectorAggNode} has no separate custom-property callback. Look for its
\texttt{BioseqScan} child and read that child's columns, pruning and filters.
The eligible shape is a plain aggregate directly above one BioSeqDB scan.
The registry covers \texttt{count(*)}, direct-column \texttt{count}, integer
\texttt{sum}, \texttt{avg}, \texttt{min}, \texttt{max}, and date
\texttt{min}/\texttt{max}. Grouping, joins, \texttt{DISTINCT}, aggregate
\texttt{FILTER}, slice projections, and aggregates over an owned search use
the ordinary aggregate route. Enabling vectorization permits the node;
it does not prove eligibility. Check NULL and empty-input results as well as
the plan shape.

\subsection{A practical reading order}

\begin{enumerate}
  \item \textbf{Literal search.} Find \texttt{BioseqScan}; check projected
    columns, pruning, vectorized filters and metadata paths. Read the index
    requirement and estimated stored/recomputed shares. Then inspect the
    scheduler, algorithms, runnable Highway target and enabled stages. With
    analysis, compare requested/actual routes, index row counts, GPU dispatches,
    declines/fallbacks, prefilter outcomes and host time. Finish with any
    column-cache snapshot.
  \item \textbf{Sequence join.} Identify each \texttt{BioseqBatchJoin} and
    its aliases. Check predicate ownership and index requirement; compare
    planned query rows/bytes with actual materialization, distinct targets and
    batch count. Read Highway and execution statistics for each node. For a
    per-row recipe, first check its explicit ownership header.
  \item \textbf{Exact cache.} Read mode and force-bypass state, parameters,
    dependencies and provider recipes. After analysis, read the actual action.
    Investigate readiness, storage, ownership and janitor status if needed.
  \item \textbf{Reusable cache.} Read branch count and actual action;
    inspect every branch action/recheck and the residual predicate. Then read
    generation state, age, TTL, rows, bytes and provider detail. A partial
    branch hit is not a full-result hit.
\end{enumerate}

For comparisons, keep two complete plans from the same data generation and
first compare the returned match rows. Then compare pipeline, architecture,
backend, index mix, prefilter work, device counters and time. A cache hit,
a different result count or a changed provider response can invalidate a
simple timing comparison. Read the PostgreSQL parents and children as part of
the same statement throughout.
